% 本文件由 LaTeX 排版 Agent 根据有效 translation.json 自动生成。
% author=Polycarp; work_id=0102; title=波利卡普的殉道

\chapter{波利卡普的殉道}

% structure: p0001 source=h1 target=omitted reason=page-title-already-rendered-by-parent

% structure: p0002 source=h2 target=section reason=relative-heading-rank; base=h2

\section{问安}

旅居斯米纳的天主的教会，致书给旅居非洛美利翁的天主的教会，并致给各地所有圣而公教会的聚会：愿从天主圣父和我们的主耶稣基督而来的怜悯、平安与爱德，广施于你们。\par

% structure: p0004 source=h2 target=section reason=relative-heading-rank; base=h2

\section{第一章 我们所写的主题}

弟兄们，我们已写信给你们，论及殉道者的事，尤其论及可敬的波利卡普，他终止了迫害，仿佛以自己的殉道为它盖上了印记。因为几乎所有先前发生的事，都是为叫主从上向我们显示一种合乎福音的殉道。他等候被交出，正如主所做的一样，好使我们也能成为他的追随者，不但顾念自己的事，也顾念他人的事。因为真诚而根基稳固的爱德之份，不但要愿意自己得救，也要愿意所有弟兄得救。\par

% structure: p0006 source=h2 target=section reason=relative-heading-rank; base=h2

\section{第二章 殉道者们奇妙的坚毅}

所有的殉道，只要是按天主的旨意发生，都是有福而高尚的。因为我们这些声称比别人更虔诚的人，理应将万有的权柄归于天主。实在，谁能不钦佩他们崇高的心志和忍耐，以及他们所表现的那份对主的爱德呢？——他们被鞭打至皮开肉绽，身体乃至内里的血管和筋络都暴露出来，却仍然忍耐，连旁边的人都怜悯哀叹他们。然而他们达到了如此崇高的气量，竟无一人发出一声叹息或呻吟；借此向我们所有人证明，基督的这些圣殉道者在遭受如此酷刑时，已超脱了身体，更准确地说，是主当时站在他们旁边，与他们交谈。并因仰望基督的恩宠，他们蔑视世间的一切痛苦，用一小时的受苦救赎了自己脱离永恒的惩罚。为此，他们野蛮行刑者的火焰在他们看来是凉爽的。因为他们心中常常记着逃避那永恒不灭之火，并以心灵之眼仰望为那忍耐之人所预备的美物：\BibleQuote{耳朵未曾听见，眼睛未曾看见，人心也未曾想到，}\BibleRef{格前 2:9}但主已向他们启示了，因为他们已不再是凡人，而已成为天使。同样，那些被判给野兽的人经受了可怕的折磨，被放在布满尖刺的床上，遭受各种其他酷刑，为的是，如果可能，暴君能通过延长的折磨使他们否认[基督]。\par

% structure: p0008 source=h2 target=section reason=relative-heading-rank; base=h2

\section{第三章 日耳曼尼库斯的坚毅。众人要求波利卡普死}

因为魔鬼确实针对他们想出了许多计谋；但感谢天主，他未能胜过所有人。因为最高贵的日耳曼尼库斯以自己的忍耐增强了别人的胆怯，并与野兽英雄般地搏斗。因为当总督试图说服他，并劝他怜惜自己的年龄时，他却主动向野兽靠拢，激怒它，渴望更快地逃离这不义和邪恶的世界。但此时，全体群众惊叹于虔诚而敬畏天主的基督徒种族所表现出的崇高心志，喊道：\Quote{除掉无神论者；找出波利卡普！}\par

% structure: p0010 source=h2 target=section reason=relative-heading-rank; base=h2

\section{第四章 背教者昆图斯}

有一个名叫昆图斯的人，是弗里吉亚人，刚从弗里吉亚来，看见野兽就害怕了。就是他强逼自己和另外几个人自愿出庭。总督多方劝说，说服了他发誓和献祭。因此，弟兄们，我们不称赞那些自愿交出自己的人，因为福音并没有这样教导。\BibleRef{玛 10:23}\par

% structure: p0012 source=h2 target=section reason=relative-heading-rank; base=h2

\section{第五章 波利卡普的离去与异象}

但最令人敬佩的波利卡普，当他初次听说（众人在搜寻他）时，毫无慌张，决定留在城内。然而，出于对许多人的尊重，他被说服离开了。于是他离开到离城不远的一处乡间别墅。他在那里与少数（朋友）在一起，日夜无所事事，只为众人和世界各地教会祈祷，这是他的惯常做法。当他祈祷时，被捉前三日，一个异象呈现在他面前；看，他头下的枕头似乎着了火。于是，他转向身边的人，先知性地对他们说：\Quote{我必定会被活活烧死。}\par

% structure: p0014 source=h2 target=section reason=relative-heading-rank; base=h2

\section{第六章 波利卡普被一个仆人出卖}

当搜寻他的人临近时，他转移到另一处住所，追赶者随即跟到了那里。他们没有找到他，就抓住了那里的两个青年，其中一个受酷刑后招认了。这样，他就不可能继续躲藏，因为出卖他的人是他自己家里的人。于是\OriginalTerm{伊勒纳尔克}{Irenarch}（其职分相当于\OriginalTerm{克勒若诺摩斯}{Cleronomus}），名叫\Person{黑落德}，急忙将他带进赛场上。[这一切都发生]为要使他完成他特殊的命运，成为基督的分享者，并且使出卖他的人遭受犹大本人的惩罚。\par

% structure: p0016 source=h2 target=section reason=relative-heading-rank; base=h2

\section{第七章 波利卡普被追捕者找到}

于是追捕他的人，连同骑兵，并带着那个年轻人，在预备日晚餐时分带着惯常的武器出发，好像出去对付一个强盗似的。\BibleRef{玛 26:55} 他们大约在傍晚到了他所在的地方，发现他正躺在一所小房子的楼上；他本可以逃到别处去，但他拒绝了，说：\Quote{愿天主的旨意成就。} \BibleRef{玛 6:10} 他听见他们来了，就下去与他们谈话。在场的人惊叹于他的高龄和坚毅，有人就说：\Quote{如此令人尊敬的老人，竟费这么大力气来捉拿吗？} 随即，就在那时，他吩咐把一些食物和饮料摆在他们面前，尽他们所需，同时恳求他们允许他安静地祈祷一个时辰。他们准许后，他站起来祈祷，充满了天主的恩宠，以至于他连续两个时辰都停不下来，使听到的人大为惊讶，以致许多人开始后悔竟然来攻击如此虔诚可敬的老人。\par

% structure: p0018 source=h2 target=section reason=relative-heading-rank; base=h2

\section{第八章 \Person{波利卡普}被带进城中}

他刚一结束祈祷，就提到了所有曾经与他有过接触的人，无论大小、尊卑，以及普世的天主教会，他的离世时刻已到，他们便把他放在一头驴上，带进城去，那日正是大安息日。治安官黑落德和他的父亲尼刻特斯（两人同乘一辆战车）遇见了他，把他接上战车，坐在他旁边，试图说服他，说：\Quote{说一句\InnerQuote{主凯撒}，并献祭，同时履行此类场合的其他仪式，以确保安全，这有什么害处呢？} 但起初他没有回答；当他们继续催逼时，他说：\Quote{我不会照你的建议做。} 他们见无法说服他，便开始对他说恶言，并粗暴地把他从战车上推下去，以致他下车时因跌倒而腿脱臼了。但他毫不慌乱，仿佛没有受苦一样，急切地全速前进，被带到竞技场，那里的喧嚣如此之大，以至于无法被人听见。\par

% structure: p0020 source=h2 target=section reason=relative-heading-rank; base=h2

\section{第九章 \Person{波利卡普}拒绝辱骂基督}

当波利卡普进入竞技场时，有一个声音从天上向他传来，说：\Quote{你要坚强，显出男子气概，波利卡普！} 没有人看见是谁对他说话；但我们在场的弟兄们听到了那声音。当他被带上前时，众人听说波利卡普被捉拿，喧嚷声更大了。他走近时，总督问他是不是波利卡普。他承认后，总督试图说服他否认基督，说：\Quote{要尊重你的年迈，} 以及按照他们的惯例，其他类似的话，如：\Quote{以凯撒的福运起誓；悔改，说：\InnerQuote{除掉无神论者。}} 但波利卡普用严峻的目光注视着竞技场中所有邪恶外邦人的群众，向他们挥手，叹息着仰望上天，说：\Quote{除掉无神论者。} 总督催逼他说：\Quote{起誓，我就释放你；辱骂基督。} 波利卡普宣告：\Quote{八十六年来我事奉祂，祂从未亏待过我；我怎能亵渎我的君王和救主？}\par

% structure: p0022 source=h2 target=section reason=relative-heading-rank; base=h2

\section{第十章 \Person{波利卡普}承认自己是基督徒}

总督再次催逼他，说：\Quote{以凯撒的福运起誓，} 他回答说：\par

\Quote{既然你徒然地催逼我，如你所说，要我以凯撒的福运起誓，并且假装不知道我是谁、是什么人，那么请听我大胆宣告：我是基督徒。你若想了解基督教教义是什么，就给我指定一天，你会听到的。}\par

总督回答说：\Quote{说服百姓。} 但波利卡普说：\par

\Quote{我认为向你陈明我的信德是合宜的；因为我们受教导，对天主所设立的权柄和掌权者，要给予一切当得的尊敬（这并不损害我们自己）。\BibleRef{罗 13:1} 至于\Emph{这些人}，我认为他们不配从我这里得到任何解释。}\par

% structure: p0027 source=h2 target=section reason=relative-heading-rank; base=h2

\section{第十一章 威胁对\Person{波利卡普}无效}

总督于是对他说：\Quote{我手头有野兽；你若不肯悔改，我就把你丢给它们。}\par

但他回答说：\Quote{那就叫它们来吧，因为我们已经不习惯弃善从恶；从恶转为义，对我是有益的。}\par

总督又对他说：\Quote{你若不肯悔改，我就用火烧死你，因为你轻视野兽。}\par

但波利卡普说：\Quote{你以那燃烧一时、片刻后即熄灭的火来威胁我，却不知道那为不敬虔之人存留的、将来的审判和永罚之火。但你为什么迟延呢？把你所预备的拿出来吧。}\par

% structure: p0032 source=h2 target=section reason=relative-heading-rank; base=h2

\section{第十二章 \Person{波利卡普}被判处火刑}

当他说这些话及许多类似的话时，心中充满自信与喜乐，容光焕发，满溢恩宠，以致不仅没有因人们对他说的话而显得忧愁，反而使总督大为惊讶，派传令官在竞技场中央三次宣告：\Quote{\Person{波利卡普}已承认他是基督徒。}传令官宣告后，住在\Place{士麦那}的全部外邦人和犹太人的群众以无法抑制的狂怒高声喊道：\Quote{这是\Place{亚细亚}的导师、基督徒之父、我们神祇的颠覆者，他一直在教导许多人不要献祭，也不要敬拜神祇。}他们这样说着，大声喊叫，恳求\Place{亚细亚}总督\Person{斐理伯}放一头狮子攻击\Person{波利卡普}。但\Person{斐理伯}回答说，他不可这样做，因为野兽表演已经结束。于是他们一致同意喊叫，应当把\Person{波利卡普}活活烧死。因为这样，那在他祈祷时关于枕头向他启示的异象必要应验：他看见枕头着火，便转过身来对与他在一起的信友预言说：\Quote{我必须被活活烧死。}\par

% structure: p0034 source=h2 target=section reason=relative-heading-rank; base=h2

\section{第十三章 火刑柴堆被竖起}

于是，这事的执行比说话更快，群众立刻从店铺和浴场收集木柴和柴捆；犹太人尤其照常热心地协助他们。当火刑柴堆准备好时，\Person{波利卡普}脱下所有衣服，松开腰带，还试图脱掉凉鞋——这是他通常不做的，因为每位信友总是争先恐后要触摸他的皮肤。因为由于他的圣洁生活，甚至在殉道之前，他已饰有一切善行。他们立刻用为火刑柴堆准备的这些材料围住他。但当他们正要把他用钉子固定时，他说：\Quote{就让我这样吧；因为赐我力量忍受烈火的那位，也会使我无需你们用钉子固定，就能在柴堆上不动摇。}\par

% structure: p0036 source=h2 target=section reason=relative-heading-rank; base=h2

\section{第十四章 \Person{波利卡普}的祈祷}

他们没有钉他，只是捆绑了他。他把双手放在背后，像一只从大羊群中取出的杰出公羊预备作祭献，并被预备为天主悦纳的全燔祭，他举目望天，说：\par

\Quote{噢，上主全能的天主，祢所钟爱的、受赞颂的圣子耶稣基督之父，藉着祂我们得以认识祢，天使与万有权能者、一切受造物、以及所有在祢面前生活的义人种族的天主，我感谢祢，因为祢视我配得上这一天、这一个时辰，得以在祢的殉道者行列中有分，在祢的基督的杯爵中，分享灵魂与身体的不朽永生，藉着圣神所赐的不朽。愿我今天在祢面前被收纳在他们当中，作为丰腴而可悦纳的祭献，正如祢，永远真实的天主，所预定的，预先向我启示的，如今又成就的。为此，我为一切事赞美祢、称颂祢、光荣祢，连同永远属天的耶稣基督，祢的爱子，与祢和圣神一起，从今世直到万世，永受光荣。\Emph{阿们}。}\par

% structure: p0039 source=h2 target=section reason=relative-heading-rank; base=h2

\section{第十五章 \Person{波利卡普}未为火所伤}

当他念完这\Emph{阿们}，结束祈祷时，受命点火的人点燃了火。当火焰猛烈喷发时，我们这些得以目睹的人看见了一个伟大的奇迹，并被保存下来以便向他人报告当时发生的事。因为火焰自行形成拱形，如同风吹满的船帆，像圆圈一样围住了殉道者的身体。他在其中显现，不像被焚烧的肉，而像被烤的面包，或像在熔炉中发光的金银。此外，我们闻到如此甜美的气味从柴堆散发，仿佛乳香或某种珍贵香料正在那里冒烟。\par

% structure: p0041 source=h2 target=section reason=relative-heading-rank; base=h2

\section{第十六章 \Person{波利卡普}被匕首刺透}

最后，当那些恶人发现他的身体不能被火烧毁时，便命令一个刽子手走近，用匕首刺透他。他这样做时，有一只鸽子飞出，并流出大量鲜血，以致火焰被熄灭；所有百姓都惊奇于不信者与蒙选者之间竟有如此差别，这位最令人钦佩的\Person{波利卡普}便是其中之一，他在我们这时代曾是宗徒性与先知性的导师，以及\Place{士麦那}天主教会的主教。因为从他口中说出的每一句话，要么已经应验，要么必将应验。\par

% structure: p0043 source=h2 target=section reason=relative-heading-rank; base=h2

\section{第十七章 基督徒被拒绝领取\Person{波利卡普}的遗体}

但当义人种族的仇敌——那嫉妒、恶毒、邪恶者——察觉到他殉道的壮观，并想到他从起初所过的无可指摘的生活，以及他如今戴着不朽的花冠、无可争议地领受了赏报，便竭力不让我们取走他丝毫的纪念物，虽然许多人都渴望这样做并拥有他的圣洁身体。为此，他怂恿\Person{尼塞特}——\Person{黑落德}的父亲、\Person{阿尔刻}的兄弟——去请求总督不要交出他的遗体埋葬，\Quote{免得，}他说，\Quote{他们离弃被钉死的那一位，转而来敬拜这个人。}他说这话是出于犹太人的建议和急切敦促，他们也监视我们，因为我们试图把他从火中取出；他们不知道这点：我们既不可能离弃那为普世得救之人而受苦的基督（无罪的为罪人），也不可能敬拜任何别的。我们确实钦崇祂为天主子；但对殉道者，作为主的门徒和追随者，我们因他们对自身君王与师傅的非凡挚爱而恰当地爱戴他们；愿我们也被作成为他们的同伴和同门徒！\par

% structure: p0045 source=h2 target=section reason=relative-heading-rank; base=h2

\section{第十八章 \Person{波利卡普}的遗体被焚}

这时，百夫长看见犹太人所激起的纷争，便将遗体置于火中焚化了。因此，我们后来拾取了他的骨骸，它们比最精美的宝石更珍贵，比黄金更纯净，我们就将之存放在一个合适的地方。在那里，当主赐予我们机会时，我们就聚集在一起，怀着喜乐与欢欣，庆祝他殉道纪念日，既为纪念那些已经跑完赛程的人，也为操练和预备那些将要步其后尘的人。\par

% structure: p0047 source=h2 target=section reason=relative-heading-rank; base=h2

\section{\Emph{第十九章　称颂殉道者波利卡普}}

以上便是对真福波利卡普的记述。他是士每拿第十二位殉道者（同时也算上非拉铁非的殉道者），然而他在众人的记忆中占有独特地位，甚至异教徒也无不称道他。他不仅是一位杰出的导师，更是一位卓越的殉道者，其殉道事迹为众人所渴望效法，因为它完全符合基督的福音。因着忍耐，他战胜了不义的官长，从而获得了不朽的荣冠。如今，他与宗徒及所有义人同在天堂，欢欣地光荣天主圣父，并赞美我们的主耶稣基督——我们灵魂的救主、我们身体的统率者，以及遍及普世天主教会的大司牧。\par

% structure: p0049 source=h2 target=section reason=relative-heading-rank; base=h2

\section{\Emph{第二十章　此书信应转达弟兄们}}

既然你们请求我们详尽告知所发生的事，我们现在通过我们的弟兄马尔谷，将这份概要寄给你们。当你们亲自阅读这封书信后，请转送给更远的弟兄们，使他们也能光荣那按自己旨意拣选仆人的主。愿那能藉祂的恩宠与仁慈，将我们众人领进祂永恒国度的主——藉祂的独生子耶稣基督——愿荣耀、尊威、权能、威严归于祂，直到永远。阿们。请向所有圣人致敬。我们这里的人问候你们，还有书写此信的厄瓦雷斯特及其全家。\par

% structure: p0051 source=h2 target=section reason=relative-heading-rank; base=h2

\section{\Emph{第二十一章　殉道日期}}

真福波利卡普是在克桑提库斯月初的第二日，即五月朔日前第七天，在大安息日第八时辰殉道的。他被黑落德逮捕，当时特拉利亚的斐理伯为大司祭，斯塔提乌斯·夸德拉图斯为总督，但耶稣基督是永远的王，愿荣耀、尊威、威严和永存的宝座归于祂，世世代代，直到永远。阿们。\par

% structure: p0053 source=h2 target=section reason=relative-heading-rank; base=h2

\section{\Emph{第二十二章　问候}}

弟兄们，愿你们一切顺利，当你们按照耶稣基督福音的教导行走时，愿荣耀藉着祂归于天主圣父和圣神，为祂神圣选民的救恩。真福波利卡普就是效法他们的榜样而受苦的，愿我们步他的后尘，也能在耶稣基督的国度中被寻得！\par

这些事是卡尤斯从依勒内（他是波利卡普的门徒）的手稿抄录的，卡尤斯本人与依勒内甚为亲密。而我，苏格拉底，在格林多从卡尤斯的手稿抄录了这些。愿恩宠与你们众人同在。\par

我，皮奥尼乌斯，又从先前写成的抄本抄录了这些，并仔细核查，真福波利卡普通过启示向我显示了它们，正如我将在下文说明。当这些事因年深日久几乎湮没时，我将它们收集起来，好使主耶稣基督也能将我与祂的选民一同收聚到祂天上的国度。愿荣耀归于祂，与圣父及圣神，直到永远。阿们。\par

\SourceRecord{Translated by Alexander Roberts and James Donaldson.；Ante-Nicene Fathers；Vol. 1.；Edited by Alexander Roberts, James Donaldson, and A. Cleveland Coxe.；Buffalo, NY: Christian Literature Publishing Co.,；1885.；<http://www.newadvent.org/fathers/0102.htm>.}
